% from main.tex
%%%%%%%%%%  scan idx 110  |  folio 97  %%%%%%%%%%
La
formule du cône \(M''{}^j=F'_i{}^{\,j+1}\oplus F^j\) montre d'autre part que \(M''\) est strictement \((i-1)\)-pseudo-cohérent. Donc, en vertu de (2.1.1), \(\Ho^{i-1}(M'')=\Ho^{i-1}(M')\) est de type fini, d'où notre assertion. On en déduit, par itération, et quitte à raffiner \(S\), un \(n\)-quasi-isomorphisme \(v_n:F'_n\longrightarrow E\), où \(F'_n\) est strictement \(n\)-pseudo-cohérent. Mais, d'après (1.4.5), on peut prolonger \(v_n\) en un quasi-isomorphisme \(v:F'\longrightarrow E\) sans rien changer en degré \(\geq n\), donc en particulier avec \(F'\) strictement \(n\)-pseudo-cohérent, ce qui prouve a).

b) L'équivalence des conditions (i) et \((\mathrm{i}')\) résulte aussitôt de a) et de la définition des isomorphismes de \(\Dpx^-(\mathcal C_X)\). Si \(E\in\Ob\Dpx^-(\mathcal C_S)\) est strictement \(n\)-pseudo-cohérent, il existe une suite exacte
\[
  0\longrightarrow E'\longrightarrow E\longrightarrow E''\longrightarrow 0,
\]
où \(E'\) est strictement parfait, et \(E''{}^i=0\) pour \(i\geq n\); par suite \(E'\to E\) est un \(n\)-quasi-isomorphisme, d'où l'implication \((\mathrm{i})\Rightarrow(\mathrm{ii})\). L'implication \((\mathrm{ii})\Rightarrow(\mathrm{ii}')\) étant triviale, il reste à prouver \((\mathrm{ii}')\Rightarrow(\mathrm{i})\). Soit donc \(u:E\longrightarrow F\) un \(n\)-isomorphisme de \(\Dpx^-(\mathcal C_S)\), avec \(E\) strictement parfait; \(u\) est donc défini par un diagramme de \(\Cpx^-(\mathcal C_S)\)
\[
\begin{tikzcd}
& E' \arrow[dl,"s"'] \arrow[dr,"u'"]\\
E && F ,
\end{tikzcd}
\]
où \(s\) est un quasi-isomorphisme et \(u'\) un \(n\)-quasi-isomorphisme. En vertu de a), il existe, quitte à raffiner \(S\), un quasi-isomorphisme \(t:E''\longrightarrow E'\), avec \(E''\) strictement \(n\)-pseudo-cohérent. On applique alors (1.4.5) pour prolonger \(u't:E''\to F\) en un quasi-isomorphisme \(E'''\longrightarrow F\) avec \(E'''\) strictement \(n\)-pseudo-cohérent, ce qui prouve \((\mathrm{ii}')\Rightarrow(\mathrm{i})\) et achève la démonstration de (2.2).

